\documentclass[10pt,a4paper]{amsart}
\usepackage[margin=19mm]{geometry}
\usepackage{amsmath,amssymb,mathtools}
\usepackage[scr=rsfs,cal=euler]{mathalfa}
\usepackage{xfrac}
\usepackage[unicode,hidelinks,bookmarks=false]{hyperref}
\newcommand{\ie}{i.e.\ }
\newcommand{\cf}{cf.\ }
\newcommand{\resp}{respectively\ }
\newcommand{\Subsection}{\subsection}
\providecommand{\nobreakdash}{\nobreak\hbox{-}}
\setlength{\emergencystretch}{2em}
\multlinegap0pt
\usepackage{bm}
\usepackage{bbm}
\usepackage{dsfont}
\usepackage{xspace}
\usepackage{cancel}
\usepackage{mathtools}
\usepackage{amsthm}
\makeatletter
\@ifundefined{definition}{\newtheorem{definition}{Definition}}{}
\@ifundefined{proposition}{\newtheorem{proposition}[definition]{Proposition}}{}
\makeatother
\newcommand{\Fq}{\mathbb F_q}
\newcommand{\ZZ}{\mathbb Z}
\newcommand{\cO}{\mathcal O}
\title[$K$-spherical horospherical averages on the Nagao quotient: ]{$K$-spherical horospherical averages on the Nagao quotient: tree combinatorics and exact discrepancy}
\author{Sanghoon Kwon}
\date{Source: arXiv:2607.08704v2 (CC BY 4.0)}
\begin{document}
\maketitle

\noindent\emph{Source and licence.} $K$-spherical horospherical averages on the Nagao quotient: tree combinatorics and exact discrepancy. Version arXiv:2607.08704v2 (CC BY 4.0), distributed under the Creative Commons Attribution 4.0 International licence, as recorded in the arXiv licence metadata for this exact version and shown on the abstract page https://arxiv.org/abs/2607.08704v2. Full rights notice: licenses/arxiv-2607.08704-CC-BY-4.0.txt.\par\smallskip

\noindent\textit{Editorial scope.} Counted unit: the Proposition whose source label is \texttt{prop:compact-normal-form} in the article (source0.tex lines 1737--1757, source label \texttt{prop:compact-normal-form}), delivered together with the article's own complete original proof of it (source0.tex lines 1759--1815, one \texttt{proof} environment of 57 lines). No mathematics is written, repaired or re-derived: statement and proof are the article's own text line for line. Required context retained uncounted: none. Every definition, display and label of those lines is retained unchanged. Omitted from this package and not redistributed: 1--1736, 1758--1758, 1816--2996. Nothing inside a retained excerpt is omitted or altered: every retained body is delivered exactly as the article's lines are written, so no omission receipt of any kind accompanies this package. Citations: the retained excerpts carry no citation command at all, so no bibliography entry of the article is transcribed and no reference is redistributed. Reference adaptation: none was needed and none was made; every reference inside the delivered unit resolves inside it, so no unresolved reference or citation is produced. Printed slips kept literal: no printed word, symbol, hypothesis or number of the counted statement or of its proof was altered. Layout and class: the article's own class is \texttt{amsart} with packages including geometry, amsmath, amssymb, amsthm, mathtools, enumitem, float, tikz, hyperref; the wrapper is the lane's pinned \texttt{amsart} editorial wrapper at 10pt on A4, which restates the theorem-like environments the retained text needs and adds the article's own macro definitions that the retained text needs (\texttt{\textbackslash Fq}, \texttt{\textbackslash ZZ}, \texttt{\textbackslash cO}), with extra wrapper packages (bm, bbm, dsfont, xspace, cancel, mathtools). The delivered numbering is the unit's own. Assets: no retained span contains a figure environment, a graphics-inclusion command, a PDF-page inclusion, a table or a TikZ picture; no image file is copied into this package, the package declares no asset dependency and no image file is redistributed with it. Licence: version arXiv:2607.08704v2 of this article is distributed under the Creative Commons Attribution 4.0 International licence, as recorded in the arXiv licence metadata for this exact version and shown on the abstract page https://arxiv.org/abs/2607.08704v2. A full rights notice is delivered at licenses/arxiv-2607.08704-CC-BY-4.0.txt. Characters: every retained excerpt is pure ASCII, so the delivered engine, XeLaTeX with its default Latin Modern text font, reports no missing character.
\begin{proposition}[Spherical normal form for compact \(U\)-orbits]
\label{prop:compact-normal-form}
Let \(Y=\Gamma gU\).  The orbit \(Y\) is compact if and only if
\(g\infty\in\mathbb P^1(\Fq(t))\).  If \(Y\) is compact, its parameter
stabilizer
\[
\Lambda_U(Y):=\{s\in F:\Gamma gu(s)=\Gamma g\}
\]
is independent of the representative and of the base point chosen on \(Y\).
It has covolume \(q^{-2m-1}\) for a unique integer \(m=m(Y)\), and for this
integer
\[
(h_K)_*\nu_Y=(h_K)_*\nu_{Y_m},
\qquad Y_m:=\Gamma a_mU.
\]
More precisely, \(m=-v_\infty(r)\), where one can write
\[
\gamma^{-1}g=d(r)u(s_0)=a_m d(\eta)u(s_0)
\]
for some \(\gamma\in\Gamma\), \(\eta\in\cO^\times\), and \(s_0\in F\).
\end{proposition}

\begin{proof}
Suppose first that \(g\infty\) is rational.  Since \(A=\Fq[t]\) is a
principal ideal domain, a Bezout identity shows that
\(\Gamma=\mathrm{SL}_2(A)\) acts transitively on
\(\mathbb P^1(\Fq(t))\).  Choose \(\gamma\in\Gamma\) with
\(\gamma\infty=g\infty\).  Then \(\gamma^{-1}g\) fixes \(\infty\), and hence
\[
\gamma^{-1}g=d(r)u(s_0)
\]
for some \(r\in F^\times\) and \(s_0\in F\).  Write \(r=t^m\eta\) with
\(m=-v_\infty(r)\in\ZZ\) and \(\eta\in\cO^\times\).  The factor \(u(s_0)\)
is absorbed by the right \(U\)-orbit, while
\[
d(\eta)u(s)=u(\eta^2s)d(\eta),
\]
so
\[
Y=\Gamma a_m d(\eta)U=Y_m d(\eta).
\]
The stabilizer of \(Y_m\) in the additive \(U\)-parameter is
\[
\Lambda_m=\{s\in F:\Gamma a_mu(s)=\Gamma a_m\}=t^{-2m}A.
\]
It is a cocompact lattice in \(F\), with fundamental domain
\(t^{-2m-1}\cO\).  Hence \(Y_m\), and therefore \(Y=Y_md(\eta)\), is compact.
Right translation by \(d(\eta)\in K\) does not change \(h_K\), and the
parameter change \(s\mapsto\eta^2s\) preserves additive Haar measure.
Consequently,
\((h_K)_*\nu_Y=(h_K)_*\nu_{Y_m}\).

The parameter stabilizer of \(Y\) is
\[
\Lambda_U(Y)=\{s\in F:gu(s)g^{-1}\in\Gamma\}
          =\eta^{-2}t^{-2m}A.
\]
Replacing \(g\) by \(\gamma_0g\), \(\gamma_0\in\Gamma\), or by
\(gu(s_1)\), \(s_1\in F\), leaves this subgroup unchanged, using
\(\Gamma\gamma_0=\Gamma\) and the commutativity of \(U\).  Its covolume under
the normalization \(ds(\cO)=1\) is
\(q^{-2m-1}\).  This covolume is intrinsic to the \(U\)-orbit and recovers
\(m\); in particular, neither the choice of \(\gamma\) nor another normal-form
presentation can change \(m(Y)\).  This proves uniqueness.

Conversely, suppose that \(Y\) is compact.  Since \(\Gamma\) is discrete,
\(\Lambda_U(Y)\) is discrete.  The orbit map
\[
\Lambda_U(Y)\backslash F\longrightarrow Y,
\qquad \Lambda_U(Y)+s\longmapsto\Gamma gu(s),
\]
is the canonical homeomorphism from the stabilizer quotient onto the closed
orbit \(Y\).  Thus compactness of \(Y\) implies that \(\Lambda_U(Y)\) is
cocompact, so it contains some \(s\ne0\).  The nontrivial unipotent matrix
\(gu(s)g^{-1}\in\mathrm{SL}_2(A)\) fixes \(g\infty\).  Its unique fixed line
is the kernel of \(gu(s)g^{-1}-I\), a matrix with entries in \(\Fq(t)\); that
line is therefore \(\Fq(t)\)-rational.  Hence
\(g\infty\in\mathbb P^1(\Fq(t))\), completing the converse.
\end{proof}

\end{document}
